% =========================================================
\subsection{Police / Community Safety}
% =========================================================
For police and community safety teams, two questions are most informative: \emph{whether the
model reads an area through its recent months or its long-run crime level}, and \emph{where
recorded crime is rising or falling} (Figure~\ref{fig:stakeholder-police}).\par\medskip
% Police / Community Safety stakeholder card (Chapter 6). The card shows the dashboard: every
% image is the dashboard's own short-card police panel for Old Town (029E), drawn by
% dashboard.views.stakeholder_figures and exported at print size by
%   python export_dashboard_panels.py --stakeholder police_short
%   01 recent-month share   02 recorded change (Old Town timeline dropped so the card fits one
%   page; its point, stable recorded crime, is in section 02). GeoSensitivity is in
%   the text and the first tile; its map is the planner card's map 01, so it is not repeated.
% Text and tiles: generated by the dashboard LLM pipeline (generate_card_narrative.py
% --stakeholder police): master prompt geosmile-hull-v8 + police card-mode addendum (v2). The
% model marks figures and key terms itself. Prompts, evidence, raw output and checks:
% police_card_llm.json (earlier run: police_card_llm_attempt1.json). The generated headline is
% not shown on the card.
%
% Preamble: \usepackage[most]{tcolorbox}, xcolor, graphicx.
% Copy into Chapter6/figs/: dashboard_police_short_E01033104_{1,2}.png and
% police_card_llm_content_openai.tex.
% =========================================================

\definecolor{cardInk}{HTML}{12171C}
\definecolor{cardMuted}{HTML}{57646D}
\definecolor{cardTeal}{HTML}{1F4E5F}
\definecolor{cardPink}{HTML}{FCE4EC}
\definecolor{cardWine}{HTML}{8A0033}
\definecolor{cardRule}{HTML}{D9DEE2}

% Same macros as planner_card.tex (\providecommand: defined once, whichever card comes first).
\providecommand{\hlpink}[1]{\begingroup\setlength{\fboxsep}{0.6pt}\colorbox{cardPink}{\textcolor{cardWine}{#1}}\endgroup}
\providecommand{\pv}[1]{\hlpink{#1}}
\providecommand{\kt}[1]{\hlpink{#1}}
\providecommand{\roletag}[1]{\tcbox[on line, colback=cardTeal, colframe=cardTeal, arc=2.5mm,
  boxsep=0pt, left=6pt, right=6pt, top=3pt, bottom=3pt]{\color{white}\sffamily\bfseries #1}}
\providecommand{\kicker}[1]{{\color{cardMuted}\sffamily\small\bfseries #1}}
\providecommand{\tile}[2]{\begin{tcolorbox}[enhanced, nobeforeafter, width=\cardtilewidth,
  equal height group=\cardtilegroup, colback=cardPink, frame hidden, arc=2mm,
  left=5pt, right=3pt, top=3pt, bottom=3pt, valign=top, before upper=\raggedright]
  {\bfseries\color{cardWine}#1}\par\vspace{1pt}
  {\sffamily\color{cardMuted}#2}\end{tcolorbox}}
% One dashboard panel under its numbered heading, image at #2 of the column width.
\providecommand{\dashpanel}[3]{\kicker{#1}\par\vspace{2pt}%
  {\centering\includegraphics[width=#2\linewidth]{#3}\par}}

\def\cardtilegroup{policetiles}
\def\cardtilewidth{0.2375\linewidth}
\providecommand{\cardHeadline}{[headline]}
\providecommand{\cardSecA}{[section 01]}
\providecommand{\cardSecB}{[section 02]}
\providecommand{\cardFooter}{[footer]}
\providecommand{\cardModel}{[model]}
\providecommand{\cardTileA}{\tile{--}{--}}
\providecommand{\cardTileB}{\tile{--}{--}}
\providecommand{\cardTileC}{\tile{--}{--}}
\providecommand{\cardTileD}{\tile{--}{--}}
\input{Chapter6/figs/police_card_llm_content_openai.tex}

\begin{figure}[!htbp]
\begin{tcolorbox}[enhanced, colback=white, colframe=cardRule, arc=3mm, boxrule=0.6pt,
                  left=3mm, right=3mm, top=2.5mm, bottom=2.5mm, before skip=0pt, after skip=2pt,
                  fontupper=\sffamily\normalsize\color{cardInk}]
\setlength{\parskip}{0pt}\raggedright%
\roletag{POLICE / COMMUNITY SAFETY}\par\smallskip
{\color{cardMuted}Sensitivity, recent or long-run signals and recorded change}\par\vspace{4pt}
% Row 1: how the model reads crime history, with section 01 beside the map.
\begin{minipage}[t]{0.46\linewidth}\vspace{0pt}
  \dashpanel{01 \textperiodcentered{} Recent or long-run signal}{0.95}{Chapter6/figs/dashboard_police_short_E01033104_1.png}
\end{minipage}\hfill
\begin{minipage}[t]{0.51\linewidth}\vspace{14pt}\raggedright
  \cardSecA\par
\end{minipage}\par\vspace{5pt}
% Row 2: what recorded crime did across Hull, with section 02 beside the map.
\begin{minipage}[t]{0.46\linewidth}\vspace{0pt}
  \dashpanel{02 \textperiodcentered{} Recorded change}{0.95}{Chapter6/figs/dashboard_police_short_E01033104_2.png}
\end{minipage}\hfill
\begin{minipage}[t]{0.51\linewidth}\vspace{14pt}\raggedright
  \cardSecB\par
\end{minipage}\par\vspace{4pt}
\kicker{Key figures}\par\vspace{2pt}
\noindent\cardTileA\hfill\cardTileB\hfill\cardTileC\hfill\cardTileD\par\vspace{4pt}
{\color{cardRule}\rule{\linewidth}{0.4pt}}\par\vspace{1pt}
{\color{cardMuted}\footnotesize \cardFooter{} Panels are the dashboard's police view for Old Town
(029E). Text, tiles and \pv{highlights} generated by \cardModel{} (prompt geosmile-hull-v8, card
mode).\par}
\end{tcolorbox}
\caption{Police / Community Safety card: the dashboard's police panels for Old Town, with the
LLM-generated narrative.}
\label{fig:stakeholder-police}
\end{figure}
